\chapter{Lumbar Puncture}

\section*{Overview}
Lumbar puncture, also known as a spinal tap, is a procedure that removes cerebrospinal fluid (CSF) from the subarachnoid space of the lumbar spine for diagnosis or to deliver medication.

\section*{Alternative Names}
Spinal tap; CSF collection; lumbar tap

\section*{Principle}
A needle is inserted between the lumbar vertebrae, typically L3-L4 or L4-L5, into the subarachnoid space. CSF pressure is measured and fluid is collected for analysis or medication such as anesthetic, chemotherapy, or antibiotics is instilled.

\section*{History}
Lumbar puncture was introduced by Heinrich Quincke in 1891 and rapidly became a cornerstone of neurologic diagnosis, especially for meningitis and subarachnoid hemorrhage.

\section*{Why It Is Done}
Performed to evaluate suspected meningitis or encephalitis, subarachnoid hemorrhage, demyelinating disease, inflammatory or neoplastic processes of the central nervous system, and to administer intrathecal medication or measure CSF pressure.

\section*{Is This a Primary or Secondary Test or Procedure}
Primary diagnostic procedure for central nervous system infection and other intrathecal pathology.

\section*{How It Is Performed}
The patient lies on the side in the fetal position or sits leaning forward. The back is prepped and anesthetized, and a spinal needle is advanced between the spinous processes until a "pop" is felt. Opening pressure is measured, CSF is collected, and the needle is removed.

\section*{Preparation}
Informed consent, review of anticoagulants and bleeding risk, and imaging before the procedure when papilledema, focal neurologic deficit, or suspected mass is present.

\section*{Contraindications and Precautions}
Increased intracranial pressure from a mass lesion, bleeding diathesis or anticoagulation, and skin infection over the puncture site are contraindications. Precaution: obtain imaging first if a space-occupying lesion is suspected.

\section*{Risks}
Post-dural-puncture headache, bleeding, infection, nerve root irritation, and rarely cerebral herniation when performed despite elevated intracranial pressure.

\section*{Aftercare and Recovery}
The patient lies flat or semireclined and is encouraged to hydrate. The sample is sent immediately for cell count, protein, glucose, culture, and other studies.

\section*{Outcome and Follow-up}
A low CSF glucose with elevated protein and neutrophils suggests bacterial meningitis; lymphocytic pleocytosis suggests viral, mycobacterial, or fungal causes; xanthochromia supports subarachnoid hemorrhage.

\section*{Expected Results}
Clear, colorless CSF with normal opening pressure, low white cell count, normal protein and glucose, and no organisms.

\section*{Follow-up}
Results guide antimicrobial and supportive therapy; repeat puncture may assess response or reclassify uncertain cases.

\section*{When to Seek Medical Care}
Follow the procedure team's specific instructions. Seek urgent medical assessment for severe or worsening pain, uncontrolled bleeding, fever or signs of infection, trouble breathing, chest pain, fainting, new weakness or confusion, inability to urinate, or any rapidly worsening symptom. The appropriate warning signs depend on the procedure and the patient's medical condition.

\section*{References}
\begin{enumerate}
  \item MedlinePlus. Lumbar puncture (spinal tap). U.S. National Library of Medicine; 2025.
  \item Current Medical Diagnosis and Treatment. McGraw-Hill; 2025.
\end{enumerate}

\clearpage
